\documentclass[11pt,a4paper]{article}

% Packages
\usepackage[utf8]{inputenc}
\usepackage[T1]{fontenc}
\usepackage[margin=0.7in]{geometry}
\usepackage{enumitem}
\usepackage{hyperref}
\usepackage{needspace}
\usepackage{titlesec}
\usepackage{xcolor}

% Formatting
\pagestyle{empty}
\setlength{\parindent}{0pt}
% Slight global indent for most bullets; skill lists keep their own small indent
\setlist[itemize]{leftmargin=12pt, noitemsep, topsep=2pt}

% Section formatting
\titleformat{\section}{\Large\bfseries}{}{0em}{}[\titlerule]
\titlespacing*{\section}{0pt}{8pt}{6pt}

\titleformat{\subsection}[runin]{\bfseries}{}{0em}{}
\titlespacing*{\subsection}{0pt}{6pt}{0pt}

% Colors - Palette professionnelle
\definecolor{primaryblue}{RGB}{0,82,155}
\definecolor{accentblue}{RGB}{0,120,215}
\definecolor{mediumgray}{RGB}{100,100,100}
% keypoint is bold only now

% Hyperlinks
\hypersetup{
    colorlinks=true,
    urlcolor=accentblue
}

% Custom commands
\newcommand{\highlight}[1]{\textcolor{primaryblue}{\textbf{#1}}}
\newcommand{\keypoint}[1]{\textbf{#1}}
% Skill title with reduced visual weight
\newcommand{\skilltitle}[1]{\normalsize\textcolor{primaryblue}{\textbf{#1}}}

\begin{document}

\begin{center}
    {\LARGE \textbf{ERIC MAYEUX}}\\
    \vspace{2pt}
    {\normalsize Scrum Master & Leader en ingénierie de la qualité}\\
    \vspace{3pt}
    43 rue Levis, Longueuil J4H 4B6, QC, Canada\\
    \href{mailto:mayeux.e@gmail.com}{mayeux.e@gmail.com} | 438-884-7645
\end{center}

\vspace{6pt}

\section*{Profil Professionnel}
\noindent
\highlight{Leader Agile} avec \keypoint{10+ ans d'expérience} combinant gestion de produit, leadership d'équipes et expertise technique approfondie en \highlight{Intégration qualité}. Atout unique : vision 360° du cycle de développement, permettant de définir des \highlight{stratégies produit réalistes}, \highlight{prioriser efficacement} et \highlight{anticiper les risques techniques}. Expert en méthodologies \highlight{Agile (Scrum, SAFe)} et alignement des \highlight{parties prenantes}. \highlight{Certifié Scrum Master}.

\Needspace{4\baselineskip}
\section*{Valeur Ajoutée : Quality Engineering comme Atout PM/SM}
\begin{itemize}[leftmargin=10pt]
    \item \highlight{Vision technique :} Priorisation réaliste du backlog en tenant compte des contraintes d'implémentation
    \item \highlight{Risques anticipés :} Critères d'acceptation testables, définition claire de la DoD
    \item \highlight{Culture DevOps :} Promeut l'intégration continue et la qualité dès la conception (shift-left)
    \item \highlight{Collaboration Dev/QA :} Facilite la communication entre équipes techniques. Adresse les enjeux de SDLC
    \item \highlight{Décisions data-driven :} KPIs qualité et tableaux de bord pour orienter la roadmap. User centric
\end{itemize}
\nopagebreak[4]
\section*{Compétences Clés}

\begin{minipage}[t]{0.55\textwidth}
\skilltitle{Gestion de produit}
\begin{itemize}[leftmargin=10pt]
    \item Priorisation et gestion du backlog; définition de la feuille de route produit (roadmap)
    \item Conception et écriture de User Story, Epic et critères d'acceptation (DoD / DoR)
    \item Analyse des besoins et comportements utilisateurs. Définition des personas
\end{itemize}

\skilltitle{Agile \& Delivery}
\begin{itemize}[leftmargin=10pt]
    \item Scrum Master certifié; facilitation de cérémonies et coaching Agile
    \item SAFe (Scaled Agile Framework), Planification d'itérations, sprints, Kanban, APOM
    \item Mise en place de pratiques TDD / Specification-Driven où pertinent
    \item Animation de pratiques. Organisation de hackathons
\end{itemize}

\skilltitle{Leadership \& Management}
\begin{itemize}[leftmargin=10pt]
    \item Management d'équipes (5+), mentorat et recrutement
    \item Animation d'ateliers, hackathons et accompagnement du changement
\end{itemize}
\end{minipage}
\hfill
\begin{minipage}[t]{0.42\textwidth}
\skilltitle{Compétences techniques}
\begin{itemize}[leftmargin=10pt]
    \item Architecture logicielle, APIs, microservices
    \item CI/CD, pipelines, containers (Docker, Kubernetes)
    \item Cloud (AWS, GCP) et stratégies de déploiement
\end{itemize}

\skilltitle{Outils \& Plateformes}
\begin{itemize}[leftmargin=10pt]
    \item Jira (cert. Admin ACP-600), Confluence, suite Atlassian
    \item GitHub, Jenkins, Datadog, Splunk
    \item Automatisation tests (Selenium, Playwright, Robot...)
\end{itemize}

\skilltitle{Parties prenantes \& Métriques}
\begin{itemize}[leftmargin=10pt]
    \item Communication client, reporting exécutif, alignment multi-équipes
    \item OKR, KPI produit, tableaux de bord et suivi de la vélocité
\end{itemize}
\end{minipage}

\section*{Expérience Professionnelle}

\subsection*{Practice Lead Qualité | Responsable Pratique Qualité \& Facilitateur Agile} \hfill \textit{\color{mediumgray}Février 2025 -- Présent}\\
Banque Nationale du Canada (BNC) -- Montréal, QC
\begin{itemize}[leftmargin=22pt]
    \item \highlight{Leadership Agile :} Pilote la pratique qualité pan-train (multi-squads) incluant équipes \textbf{internationales (Thaïlande)}
    \item \highlight{Management :} Encadre, mentore et développe une équipe de \keypoint{5+ professionnels QA} en environnement SAFe
    \item \highlight{Stratégie produit :} Définit et exécute la stratégie de test alignée avec la feuille de route produit
    \item \highlight{Facilitation :} Anime des ateliers de définition de la \textbf{DoD (Definition of Done)} et des critères d'acceptation
    \item \highlight{Décisions basées sur les données :} Implémente des KPI qualité et tableaux de bord (Datadog, Splunk, Jira)
    \item \highlight{Innovation produit :} Conduit des POC (\textbf{IA, Playwright MCP, Crew AI}) et évalue l'impact business
    \item \highlight{Gestion du changement :} Organise un hackathon de test et promeut une culture d'amélioration continue
    \item \highlight{DevOps :} Met en place des pipelines CI/CD/CT améliorant le \keypoint{time-to-market de 40\%}
\end{itemize}

\subsection*{Intégrateur Qualité Senior | Ingénieur Qualité Senior \& Chef de projet technique} \hfill \textit{\color{mediumgray}Janvier 2022 -- Février 2025}\\
Banque Nationale du Canada (BNC) via Groupe Astek -- Montréal, QC
\begin{itemize}[leftmargin=22pt]
    \item \highlight{Collaboration PO :} Travaille étroitement avec les Product Owners pour définir des User Stories testables
    \item \highlight{Affinage du backlog :} Participe au grooming et apporte une perspective technique pour une priorisation efficace
    \item \highlight{Stratégie produit :} Conçoit la stratégie de test alignée avec la feuille de route Gestion de Patrimoine \& ECM
    \item \highlight{Gestion des parties prenantes :} Présente les métriques qualité aux parties prenantes pour assurer la transparence sur l'état du produit
    \item \highlight{Gestion des risques :} Identifie les risques techniques dès la conception, évitant des retards de livraison
    \item \highlight{Expertise technique :} Conseille le PO sur la faisabilité technique et l'effort estimé
    \item \highlight{Reporting :} Crée des tableaux de bord qualité permettant des décisions basées sur les données
\end{itemize}

\subsection*{QA Lead | Responsable Automatisation QA} \hfill \textit{\color{mediumgray}Juin 2021 -- Décembre 2021}\\
AODocs -- Paris, France
\begin{itemize}[leftmargin=22pt]
    \item \highlight{Optimisation de la livraison :} Automatise \keypoint{100\% des tests de régression}, libérant \keypoint{20h/sprint} pour de nouvelles fonctionnalités
    \item \highlight{POC et innovation :} Évalue de nouveaux outils (tests mobiles Flutter, API Karate) et recommande leur adoption
    \item \highlight{Standards qualité :} Implémente des design patterns (Page Object Model) adoptés par toute l'équipe
    \item \highlight{Transparence :} Génère des rapports Allure facilitant la communication sur l'état du produit aux parties prenantes
\end{itemize}

\subsection*{Ingénieur DevOps \& Coordinateur de projet technique} \hfill \textit{\color{mediumgray}Janvier 2020 -- Juin 2021}\\
ENGIE Digital via Groupe Hénix -- Paris, France
\begin{itemize}[leftmargin=22pt]
    \item \highlight{Gestion de projet :} Administre Jira (\keypoint{8 projets, 50+ utilisateurs}), pilote la feuille de route et priorise les initiatives
    \item \highlight{KPI et métriques :} Met en place des tableaux de bord pour suivre la vélocité, les bugs et les performances des pipelines
    \item \highlight{Stratégie CI/CD :} Conçoit une chaîne complète (Jenkins, Docker, Kubernetes) réduisant le \keypoint{time-to-market de 60\%}
    \item \highlight{Évaluation d'outils :} Réalise des POC Robot Framework et Xray, recommande l'adoption basée sur le \highlight{retour sur investissement (ROI)}
    \item \highlight{Collaboration transverse :} Coordonne les équipes Dev, Ops et QA pour une intégration continue
\end{itemize}

\subsection*{Product Owner \& Responsable Delivery} \hfill \textit{\color{mediumgray}Juin 2019 -- Septembre 2019}\\
Hénix -- Montrouge, France
\begin{itemize}[leftmargin=22pt]
    \item \highlight{Product Owner :} Gère le backlog produit et priorise les fonctionnalités du plugin Xsquash (Jira avec Squash TM)
    \item \highlight{Histoires utilisateur :} Rédige et valide les histoires utilisateur pour l'intégration bi-directionnelle Jira/Squash
    \item \highlight{Succès client :} Assure le support SaaS et collecte les retours utilisateurs pour une amélioration continue
    \item \highlight{Migration à grande échelle :} Participe à la reprise du Centre de Service Thales (\keypoint{8500 utilisateurs}, 200+ projets)
\end{itemize}

\subsection*{Développeur logiciel | Développeur Full-Stack} \hfill \textit{\color{mediumgray}Mai 2017 -- Mai 2019}\\
MEDIAPOST via Groupe Hénix -- Paris, France
\begin{itemize}[leftmargin=22pt]
    \item Développe des \highlight{webservices SOAP/REST} (\highlight{SOA BPEL}) pour des applications de \highlight{contractualisation} et de \highlight{facturation}
    \item Crée une \highlight{application web} de gestion comptable (\highlight{Node.js}) avec \highlight{import/export} multi-formats (\highlight{CSV, XML})
    \item Automatise des \highlight{processus métier} avec \highlight{Google Apps Script} et macros Excel (gain \keypoint{15h/semaine})
    \item Collabore avec l'équipe produit pour définir des \highlight{spécifications techniques} et des \highlight{parcours utilisateur}
\end{itemize}

\subsection*{Administrateur Plateformes \& Outils | Coordinateur technique} \hfill \textit{\color{mediumgray}Avril 2016 -- Avril 2017}\\
ENGIE IT via Groupe Hénix -- Saint-Ouen, France
\begin{itemize}[leftmargin=22pt]
    \item Gère les \highlight{environnements de production} et coordonne les \highlight{déploiements} avec des équipes multiples
    \item Administre \highlight{HP ALM Performance Center} (\keypoint{150+ utilisateurs}) et assure le \highlight{support technique}
    \item Pilote les \highlight{évolutions de la plateforme} basées sur les besoins utilisateurs et le \highlight{retour d'expérience}
\end{itemize}
\vspace{6pt}

\section*{Certifications \& Formations}

\textbf{PMP (en cours)} \hfill \textit{\color{mediumgray}2025}\\
\textbf{Scrum Fundamentals Certified (SFC)} \hfill \textit{\color{mediumgray}2024}\\
\textbf{Jira ACP-600 (Administration Jira)} \hfill \textit{\color{mediumgray}2019}\\
Autres : NeoLoad Pro, Selenium Professional, SoapUI, RPA \hfill \textit{\color{mediumgray}2019}

\section*{Formation Académique}

\textbf{Cursus Automaticien et DevOps} \hfill \textit{\color{mediumgray}2019}\\
AFCEPF, Montrouge, France -- Formation DevOps, CI/CD, Infrastructure as Code. 

\textbf{Master II (niveau) -- Architecture logicielle \& Analyste informaticien} \hfill \textit{\color{mediumgray}2017}\\
AFCEPF, France -- UML, Bases de données, Architectures distribuées, Big Data

\textbf{Licence en Gestion | Bachelor's Degree in Management} \hfill \textit{\color{mediumgray}2011}\\
École Supérieure de Gestion (E.S.G) -- Paris, France

\vspace{6pt}

\section*{Compétences Linguistiques}

\textbf{Français :} Langue maternelle \quad | \quad \textbf{Anglais :} Niveau avancé (Compétences professionnelles)

\end{document}